\scanpage{105}
% Archive foliation: 106. Reused unrelated typescript; printed page 19, n°382.
\begin{center}-- 19 --\hfill n\textsuperscript{o} 382\end{center}
de $a$ ; si $Z(f,g)$ désigne l'ensemble des points vérifiant ces conditions, $Z(f,g)$ est \uline{ouvert et fermé} dans $U$ ; en particulier, si $U$ est connexe, $Z(f,g)$ est soit vide, soit égal à $U$ (``principe du prolongement analytique'').

3.6 -- \uline{Critères d'analyticité}.

3.6.1. Supposons que $K=\underline{\underline{R}}$. Soit $U$ un ouvert de $\underline{\underline{R}}^n$ et soit $f$ une application de $U$ dans $E$. Les deux conditions suivantes sont équivalentes :

(1) $f$ est \uline{analytique} dans $U$.

(2) $f$ est \uline{indéfiniment différentiable} dans $U$ (cf. n\textsuperscript{o} 2.3.1) et pour toute partie compacte $F$ de $U$, il existe des nombres réels $C$ et $M$ tels que
\[
\lVert\Delta^s f(x)\rVert\leq C.M^{|s|}\qquad\text{si }x\in F,\quad s\in\underline{\underline{N}}^n.
\]

3.6.2. Supposons que $K=\underline{\underline{C}}$. Toute application \uline{dérivable} d'un ouvert $U$ de $\underline{\underline{C}}^n$ dans $E$ est \uline{analytique} dans $U$.

\begin{center}-:-:-:-\end{center}
